\section{Cancellation for every prescribed positive scale}

Let
\[
 \Theta=\{0,-1/2,-1/4,-2/5,-1,-1/3\},\qquad
 F(t)=\frac{t(t+1)(2t+1)(3t+1)(4t+1)(5t+2)}{120},
 \qquad R(t)=\frac{F(t)}{(t-1)^8}.
\]
Define the rational involution
\[
                         \psi(t)=-\frac{t}{5t+1}.
\]
It fixes $0,-2/5$ and exchanges $-1/2$ with $-1/3$ and
$-1/4$ with $-1$.

\begin{lemma}
At every rational argument where the expressions are defined,
\[
 F(\psi(t))=-\frac{F(t)}{(5t+1)^6},\qquad
 R(t)R(\psi(t))
 =-\left(\frac{F(t)(5t+1)}
               {(t-1)^4(6t+1)^4}\right)^2.
\]
\end{lemma}
\begin{proof}
Put $N=-t$ and $D=5t+1$. Then
\[
 N+D=4t+1,\quad 2N+D=3t+1,\quad 3N+D=2t+1,\quad
 4N+D=t+1,\quad 5N+2D=5t+2.
\]
Thus the product of $N$ and these five factors is $-120F(t)$.
Division by $120D^6$ gives the first formula. Since
$\psi(t)-1=-(6t+1)/(5t+1)$, the second follows immediately.
\end{proof}

\begin{theorem}
For every rational $a>4$, put
\[
                     \lambda=-\frac{a+1}{5a}.
\]
Then $\lambda,a\lambda$ avoid every contact and the pole of $R$,
$a\lambda=\psi(\lambda)$, and $R(\lambda)R(a\lambda)$ is the
negative of a nonzero rational square. In particular every finite
valuation of this product is even. More explicitly,
\[
 R(\lambda)R(a\lambda)=
 -\left(
 \frac{25a(a-4)(a-1)(a+1)(2a-3)(3a-2)(4a-1)}
      {24(a+6)^4(6a+1)^4}
 \right)^2.
\]
\end{theorem}
\begin{proof}
One has $5\lambda+1=-1/a$ and
$\psi(\lambda)=-(a+1)/5=a\lambda$.
The inequalities
\[
                  -1/4<\lambda<-1/5,\qquad a\lambda<-1
\]
exclude all six contacts. They also exclude the pole $1$.
The transformed denominators do not vanish:
\[
 \lambda-1=-\frac{6a+1}{5a},\qquad
 6\lambda+1=-\frac{a+6}{5a}.
\]
The lemma proves the negative-square assertion. Substitution and
cancellation give the displayed square root, whose factors are all
nonzero when $a>4$.
\end{proof}

Consequently a restriction to the scales in the weak-$abc$ source,
\[
                         a=\beta^8+T(n)^8,
\]
does not prevent this cancellation. Indeed $\beta\ge3$ and $T(n)>6$
make $a>4$, and the theorem gives an admissible rational slope for
every such scale. Fixing any nonzero index $n$ and allowing the
rational input to tend to infinity gives scales of unbounded height
whose product odd radical remains one. Thus no height inequality
\[
 h(a)\le C\log\operatorname{Odd}(R(\lambda)R(a\lambda))+C_0
\]
can hold even on this prescribed-scale domain. The conclusion is
unchanged after multiplying the product by a rational square or
omitting any fixed finite set of primes from the odd radical.

This is an obstruction to the stated product-based height argument.
It does not assert that the remaining elliptic test equations or the
pole predicate have witnesses on this family. In particular the
product is negative, whereas the completeness construction uses
positive pole inputs.
